\bookchapter{第七章：最先端機能の完全解析}

\chapterintro{対応機種：}

\booksection{Apple}

：iPhone 14シリーズ以降、衛星経由での緊急SOS送信をサポート。iPhone 16シリーズはさらに衛星経由でのiMessageとSMS送信をサポート。

\booksubsection{対応機種：}

：Mate 50シリーズから北斗衛星メッセージ（テキストのみ）をサポート、Mate 60 Proシリーズはさらに天通衛星通話（電話可能）をサポート、Mate 70シリーズは衛星通信能力がさらに強化。

\booksubsection{注意事項：}

\begin{itemize}
  \item \textbf{Apple}：iPhone 14シリーズ以降、衛星経由での緊急SOS送信に対応しています。iPhone 17シリーズはさらに衛星経由のiMessageとSMS送信に対応し、より多くの地域でのカバレッジを拡大しています。
  \item \textbf{Huawei}使用シーン：
  \item \textbf{アウトドア探検、登山中の遭難。}海上での遭難。
  \item \textbf{自然災害による地上通信の中断。}使い方：
\end{itemize}

\booksubsection{開けた場所にいることを確認し、空が見えること。}

\begin{itemize}
  \item スマートフォンの衛星通信機能をオンにする（通常は緊急発信インターフェースにある）。
  \item 画面の指示に従い、衛星の方向に向ける。
  \item 事前設定された救助要請メッセージと位置を送信。
\end{itemize}

\booksubsection{豆知識}

\begin{enumerate}
  \item 衛星通信にはスマートフォンに特殊なハードウェアアンテナが必要で、すべてのスマートフォンがサポートしているわけではない。また、衛星通信は通常有料のサブスクリプションサービスが必要だが、緊急救助機能は一般的に無料。
  \item 二、ネットワーク不要通信：近距離通信の新方式
  \item 「ネットワーク不要通信」とは、モバイルネットワークやWi-Fiがない状況で、2台のスマートフォン間で直接通信する技術を指す。
  \item 技術原理：
\end{enumerate}

\begin{tcbinfo}{Bluetoothや専用の無線周波数を利用し、近距離（通常数百メートルから数キロメートル）でポイントツーポイント接続を確立。}
応用シーン：

\end{tcbinfo}

\booksection{アウトドア活動}

：ハイキング、キャンプ時のチームメンバー間の通信。

\booksubsection{災害救援}

：地震、洪水後の通信中断時。

\booksubsection{活用シーン：}

\begin{itemize}
  \item \textbf{：コンサート、スポーツイベントで人が多く信号が悪い時。}対応プラットフォーム：
  \item \textbf{Apple}：iOS 17の「NameDrop」と「オフラインで友達を探す」機能。
  \item \textbf{大型イベント}：Changlianアプリがネットワーク不要通話をサポート。
\end{itemize}

\booksubsection{対応プラットフォーム：}

\begin{itemize}
  \item \textbf{Apple}三、衝突検出：決定的瞬間に命を救う
  \item \textbf{Huawei}動作原理：
  \item \textbf{Xiaomi}カウントダウンインターフェースをポップアップし、あなたが無事かどうかを確認。
\end{itemize}

\booksection{カウントダウン終了時に応答がない場合、自動的に緊急電話を発信。}

緊急連絡先にあなたの位置を送信。

\booksubsection{対応機種：}

Apple

\begin{enumerate}
  \item ：iPhone 14シリーズ以降。
  \item カウントダウン内に応答がない場合、自動で緊急電話を発信します。
  \item ：一部のハイエンド機種が順次サポート。
\end{enumerate}

\booksubsection{注意事項：}

\begin{itemize}
  \item \textbf{Apple}誤検出の可能性あり（ジェットコースター、激しい運動など）、その際は速やかにキャンセルすればよい。
  \item \textbf{Android}四、星閃技術：次世代ワイヤレス接続
\end{itemize}

\booksubsection{星閃（SparkLink）はHuaweiが主導して開発した次世代ワイヤレス通信技術で、特定のシーンでBluetoothとWi-Fiを代替することを目指す。主要指標で既存のソリューションに匹敵、あるいは超越：空中速度は900Mbpsに達し（Bluetoothの50Mbpsを遥かに超越）、レイテンシは20 $\mu$ sまで低減（Bluetoothは約15ms）、単一ネットワークで4096デバイスの同時接続をサポート（Bluetoothは通常8台以内）。}

\begin{itemize}
  \item 技術的特徴：
  \item 超低レイテンシ
  \item ：エンドツーエンドのレイテンシはわずか20マイクロ秒（0.02ms）、Bluetoothより2桁低く、ゲームや音楽の同期はほぼ知覚できない。
\end{itemize}

\booksection{高速}

：単一リンクのピーク速度900Mbps、Bluetoothオーディオの数十倍で、ロスレスオーディオやHD動画の転送が可能。

\booksubsection{低消費電力}

\begin{itemize}
  \item \textbf{：待機時の消費電力はBluetoothと同等だが、伝送効率が高く、単位データあたりのエネルギー消費が低い。}広カバレッジ+高同時接続
  \item \textbf{：接続距離はBluetoothより50\%以上長く、単一ネットワークで数千デバイスの同時オンラインをサポート。}応用シーン：
  \item \textbf{ワイヤレスイヤホン}：より低いレイテンシ、ゲーム時に音画同期。
  \item \textbf{スマートウォッチ}：より安定した接続。
\end{itemize}

\booksubsection{活用シーン：}

\begin{itemize}
  \item \textbf{：より多くのデバイスの同時接続。}車載連携
  \item \textbf{：よりスムーズなスマートフォン-カーナビ接続。}対応デバイス：
  \item \textbf{スマートホーム}技術展望
  \item \textbf{星閃技術はワイヤレス通信の未来の方向を代表する。より多くのメーカーが参加することで、星閃対応デバイスが増え、より優れたワイヤレス体験を享受できるようになるだろう。}確認してみよう：あなたのスマートフォンがこれらの機能をサポートしているか
\end{itemize}

\booksubsection{設定で「衛星」を検索し、衛星通信関連のオプションがあるか確認。}

「衝突検出」や「緊急SOS」を検索し、有効になっているか確認。

\textbf{スマートフォンがHuaweiの場合、星閃デバイス接続をサポートしているか確認。}\par

これらの機能の使用方法を理解し、いざという時に役立てる。

\begin{tcbpractice}{二、ソフトウェア編（第8-14章）}
\begin{enumerate}
  \item 基本操作をマスターした後、各種アプリの使用テクニックを深く理解しよう。通信からオフィス、決済からエンターテインメントまで、これらの章はあなたを真のスマートフォン達人にする。
  \item 第八章：通信：SMS、電話とソーシャル
  \item 通信はスマートフォンの最も原始的で最も核心的な使命だ。現在、それは単純な通話とSMSから、インスタントメッセージ、ソーシャルネットワーク、電子メールを含む複雑なエコシステムへと進化した。
  \item 一、電話アプリの高度な機能
\end{enumerate}

\end{tcbpractice}
